\documentclass[10pt]{beamer}
\usepackage{tikz}
\usepackage{pgfplots}
\usepackage{booktabs}
\usepackage{multirow}
\usepackage{array}
\usepackage{siunitx}
\usepackage{makecell}
\usepackage{CJKutf8}

\pgfplotsset{compat=newest}
\setbeamertemplate{navigation symbols}{}

% Add section pages
\AtBeginSection[]{
  \begin{frame}
    \vfill
    \centering
    \begin{beamercolorbox}[sep=8pt,center,shadow=true,rounded=true]{title}
    \usebeamerfont{title}\insertsectionhead\par%
    \end{beamercolorbox}
    \vfill
  \end{frame}
}

% Add page number to the lower right corner
\setbeamertemplate{footline}{%
  \hfill\usebeamercolor[fg]{page number in head/foot}%
  \usebeamerfont{page number in head/foot}%
  \insertframenumber\,/\,\inserttotalframenumber\hspace*{1ex}\vskip2pt%
}

% Customize the page number appearance
\setbeamercolor{page number in head/foot}{fg=gray}
\setbeamerfont{page number in head/foot}{size=\small}

\title{Programming Languages \& Tools in Data Science / Economics}
\author{Tzu-You Lin}
\setbeamerfont{title}{size=\huge}
\date{April 16, 2025}

\begin{document}

\begin{frame}
    \titlepage
\end{frame}

\begin{frame}
    \frametitle{Outline}
    \tableofcontents
\end{frame}

\section{Python}

\begin{frame}
    \frametitle{Python: General-Purpose Language}
    \begin{itemize}
        \item Most popular programming language in 2025
        \item Easy to learn with clean, readable syntax
        \item Industry standard for modern data analysis:
            \begin{itemize}
                \item Widely adopted by major corporations
                \item Extensive data science ecosystem
                \item Large community support
            \end{itemize}
    \end{itemize}
\end{frame}

\begin{frame}
    \frametitle{Python: Ecosystem \& Libraries - Data Manipulation}
    \begin{itemize}
        \item \texttt{pandas}
            \begin{itemize}
                \item Powerful data structures for handling structured data
                \item Efficient data manipulation and analysis tools
            \end{itemize}
        \item \texttt{numpy}
            \begin{itemize}
                \item Fast numerical operations and array computations
                \item Foundation for many scientific computing libraries
            \end{itemize}
    \end{itemize}
\end{frame}

\begin{frame}
    \frametitle{Python: Ecosystem \& Libraries - Machine Learning}
    \begin{itemize}
        \item \texttt{scikit-learn}
            \begin{itemize}
                \item Comprehensive machine learning toolkit
                \item Simple and efficient tools for data analysis
                \item Supports many machine learning algorithms out of the box
            \end{itemize}
        \item \texttt{TensorFlow}
            \begin{itemize}
                \item Industry standard for deep learning
                \item Robust production deployment capabilities
            \end{itemize}
        \item \texttt{PyTorch}
            \begin{itemize}
                \item Favored in academia for research
                \item Dynamic computational graphs
                \item Machine learning papers are more likely to use PyTorch
            \end{itemize}
    \end{itemize}
\end{frame}

\begin{frame}
    \frametitle{Python: Ecosystem \& Libraries - Visualization}
    \begin{itemize}
        \item \texttt{matplotlib}
            \begin{itemize}
                \item Foundation for Python plotting
                \item Highly customizable static visualizations
            \end{itemize}
        \item \texttt{seaborn}
            \begin{itemize}
                \item Statistical data visualization
                \item Built on matplotlib with higher-level interface
            \end{itemize}
        \item \texttt{plotly}
            \begin{itemize}
                \item Interactive and web-based visualizations
                \item Modern, publication-quality graphs
            \end{itemize}
    \end{itemize}
\end{frame}

\begin{frame}
    \frametitle{Python: Industry Adoption}
    \begin{itemize}
        \item \textbf{Community Support}
            \begin{itemize}
                \item Extensive documentation and tutorials
                \item Large community of developers and users
                \item Vast number of open-source packages
            \end{itemize}
        \item \textbf{Development Environment}
            \begin{itemize}
                \item Interactive environments like Jupyter Notebook
                \item Ideal for rapid prototyping
                \item Easy to share and collaborate
            \end{itemize}
    \end{itemize}
\end{frame}

\begin{frame}
    \frametitle{Python: What should I learn next?}
    \begin{itemize}
        \item \textbf{Web Scraping \& Regular Expressions:}
            \begin{itemize}
                \item Learn \texttt{requests} and \texttt{BeautifulSoup} for web scraping
                \item Master regular expressions for text pattern matching
                \item Practice with \texttt{pandas} read\_html for structured web data
            \end{itemize}
        \item \textbf{Advanced Data Visualization:}
            \begin{itemize}
                \item Deep dive into \texttt{matplotlib} customization
                \item Master \texttt{seaborn} statistical plotting
                \item Learn to create publication-ready figures
            \end{itemize}
        \item \textbf{Deep Learning:}
            \begin{itemize}
                \item Follow Dr. Hung-yi Lee's Machine Learning course
                \item Start with neural network fundamentals
                \item Progress to advanced topics like CNN, RNN, and transformers
            \end{itemize}
    \end{itemize}
\end{frame}

\section{R}

\begin{frame}
    \frametitle{What is R?}
    \begin{itemize}
        \item \textbf{Programming Language for Statistics:}
            \begin{itemize}
                \item Open-source language developed by statisticians for statisticians
                \item Specialized in statistical computing and data analysis
                \item Strong focus on data manipulation and visualization
                \item It is the go-to language for statistics for about 10 years.
            \end{itemize}
        \item \textbf{Key Strengths:}
            \begin{itemize}
                \item Interactive data analysis environment
                \item Extensive statistical and mathematical functions
                \item Powerful graphics capabilities
                \item Large collection of specialized packages (CRAN)
            \end{itemize}
        \item \textbf{Academic and Research Focus:}
            \begin{itemize}
                \item Widely used in academic research
                \item Strong presence in statistics and economics
                \item Excellent for reproducible research
            \end{itemize}
    \end{itemize}
    This would be the most-used language in NTU ECON, you should be learning it in your second year.
\end{frame}

\begin{frame}
    \frametitle{R}
    \begin{itemize}
        \item \textbf{Popular Packages and Tools:}
            \begin{itemize}
                \item \texttt{ggplot2} for layered, grammar-of-graphics-based visualization.
                \item \texttt{tidyverse} suite (including \texttt{dplyr}, \texttt{tidyr}, \texttt{readr}) for streamlined data manipulation and cleaning.
                \item \texttt{knitr} and \texttt{rmarkdown} for dynamic reporting and reproducible research documents.
                \item Specialized packages for econometrics (\texttt{lm}, \texttt{lme4}, \texttt{forecast}).
            \end{itemize}
        \item \textbf{Industry Applications:}
            \begin{itemize}
                \item Growing adoption in data science and business analytics
                \item Integration capabilities with other programming languages
                \item Robust support for big data processing
            \end{itemize}
        \item \textbf{Learning Resources:}
            \begin{itemize}
                \item Rich online community and documentation
                \item Active Stack Overflow presence
                \item Many free tutorials and courses available
            \end{itemize}
    \end{itemize}
\end{frame}

\begin{frame}
    \frametitle{R: What to Learn?}
    \begin{itemize}
        \item \textbf{Prerequisites:}
            \begin{itemize}
                \item Basic statistical knowledge is helpful
                \item Understanding of data analysis concepts
            \end{itemize}
        \item \textbf{Learning Path:}
            \begin{itemize}
                \item Start with base R syntax and fundamental data structures
                \item Master data frame operations and manipulation
                \item Learn essential packages:
                    \begin{itemize}
                        \item \texttt{tidyverse} for modern data manipulation
                        \item \texttt{ggplot2} for advanced data visualization
                    \end{itemize}
                \item Progress to statistical modeling and regression analysis
            \end{itemize}
        \item \textbf{Career Relevance:}
            \begin{itemize}
                \item Particularly valuable for economics/econometrics research
                \item Strong foundation for academic research methodology
                \item Essential for quantitative analysis in economics
            \end{itemize}
    \end{itemize}

    Following Dr. Shiu-Sheng Chen's {\small \begin{CJK*}{UTF8}{bsmi}（陳旭昇教授）\end{CJK*}} book: {\small \begin{CJK*}{UTF8}{bsmi}機率與統計推論: R語言的應用\end{CJK*}} might be a good idea.
\end{frame}


\section{Stata}

\begin{frame}
    \frametitle{Stata}
    \begin{itemize}
        \item \textbf{Specialized Software for Econometrics:}
            \begin{itemize}
                \item Tailored for statistical analysis with a focus on reproducibility.
                \item Integrated environment designed for econometric research.
            \end{itemize}
        \item \textbf{Key Features:}
            \begin{itemize}
                \item Comprehensive support for panel data analysis, time-series analysis, and causal inference.
                \item Robust inbuilt commands with well-documented procedures.
                \item High reliability and consistent output, ideal for academic research.
                \item Usually has the latest econometric methods' implementation.
            \end{itemize}
        \item \textbf{Considerations:}
            \begin{itemize}
                \item Proprietary software with high licensing costs.
                \item Limited flexibility and customizability compared to open-source alternatives.
            \end{itemize}
    \end{itemize}
\end{frame}

\begin{frame}
    \frametitle{Should I Learn Stata?}
    \begin{itemize}
        \item \textbf{Consider Learning Stata If:}
            \begin{itemize}
                \item Your courses specifically require Stata usage
                \item You have the need to use the latest econometric methods out of the box
                \item Your research team or advisor primarily uses Stata
            \end{itemize}
        \item \textbf{Otherwise, R May Be Sufficient:}
            \begin{itemize}
                \item R provides similar statistical capabilities
                \item Open-source with no licensing costs
                \item Larger community and more extensive package ecosystem
                \item Better integration with modern data science workflows
            \end{itemize}
        \item \textbf{Decision Factors:}
            \begin{itemize}
                \item Evaluate your specific research requirements
                \item Consider the cost-benefit of learning multiple tools
                \item Assess the prevalent tools in your target research area
            \end{itemize}
    \end{itemize}
\end{frame}

\section{SQL}

\begin{frame}
    \frametitle{SQL}
    \begin{itemize}
        \item \textbf{Database Querying Language:}
            \begin{itemize}
                \item Essential for interacting with relational databases.
                \item Standardized language across various systems like MySQL, PostgreSQL, Oracle, and SQL Server.
            \end{itemize}
        \item \textbf{Core Concepts:}
            \begin{itemize}
                \item Data Retrieval: Using SELECT, FROM, and WHERE to filter and extract data.
                \item Data Aggregation: GROUP BY and HAVING for summarizing information.
                \item Data Modification: Commands for INSERT, UPDATE, DELETE operations.
                \item Advanced Operations: JOINs for combining tables, subqueries, and window functions for complex queries.
            \end{itemize}
        \item \textbf{Best Practices:}
            \begin{itemize}
                \item Optimize queries through indexing and query tuning.
                \item Normalize data to minimize redundancy.
                \item Use transactions to ensure data integrity.
            \end{itemize}
    \end{itemize}
\end{frame}

\begin{frame}
    \frametitle{Should I Learn SQL?}
    \begin{itemize}
        \item \textbf{Career Perspective:}
            \begin{itemize}
                \item Essential requirement for most Data Science positions
                \item Standard skill requirement for data-related internships
                \item Fundamental knowledge for corporate data analysis roles
            \end{itemize}
        \item \textbf{When to Learn SQL:}
            \begin{itemize}
                \item Planning to pursue a career in Data Science
                \item Interested in corporate data analysis positions
                \item Need to work with large-scale structured data
            \end{itemize}
        \item \textbf{Learning Investment:}
            \begin{itemize}
                \item Relatively straightforward syntax compared to programming languages
                \item Many free online resources available
                \item Skills remain relevant across different database systems
            \end{itemize}
        \item \textbf{Practical Value:}
            \begin{itemize}
                \item Efficient data manipulation and analysis
                \item Complements other data science tools like Python and R
            \end{itemize}
    \end{itemize}
\end{frame}

\begin{frame}
    \frametitle{How to Learn SQL?}
    \begin{itemize}
        \item \textbf{LeetCode SQL 50:}
            \begin{itemize}
                \item Curated collection of essential SQL problems
                \item Organized by difficulty level
                \item Interactive practice with immediate feedback
                \item Great starting point for beginners
            \end{itemize}
    \end{itemize}
\end{frame}




\section{PowerBI and Tableau}

\begin{frame}
    \frametitle{PowerBI and Tableau}
    \begin{itemize}
        \item \textbf{Business Intelligence Tools:}
            \begin{itemize}
                \item Designed for interactive data visualization and dashboard creation.
                \item Widely used in corporate settings for real-time analytics.
            \end{itemize}
        \item \textbf{Key Capabilities:}
            \begin{itemize}
                \item Drag-and-drop interfaces that simplify dashboard design.
                \item Seamless integration with diverse data sources, including cloud services, databases, and spreadsheets.
                \item Advanced features such as predictive analytics, trend analysis, and drill-down capabilities.
            \end{itemize}
        \item \textbf{Comparative Highlights:}
            \begin{itemize}
                \item \textbf{PowerBI:} Strong integration with the Microsoft ecosystem, cost-effective licensing, and frequent updates.
                \item \textbf{Tableau:} Emphasizes visual storytelling with flexible data connections and a robust user community.
            \end{itemize}
        \item \textbf{Use Cases:}
            \begin{itemize}
                \item Real-time monitoring of key performance indicators (KPIs).
                \item Creating tailored dashboards for diverse business needs.
                \item Facilitating collaborative data exploration and decision-making.
            \end{itemize}
    \end{itemize}
\end{frame}

% \begin{frame}
%     \frametitle{References}
% \end{frame}

\end{document}
